\begin{proofbox}
\textbf{\textcolor{CProof}{统一记号}}
令 $m=|PF_1|>0$、$n=|PF_2|>0$、$|F_1F_2|=2c$，$\theta=\angle F_1PF_2$，$S=S_{\triangle F_1PF_2}$。始终有
\[
|m-n|=2a,\qquad c^2=a^2+b^2,\qquad 0<\theta<\pi.
\]
以下五种证明从不同角度组织思路；\textbf{考场计算首选方法一，零余弦定理的几何理解首选方法五，规范设点首选方法三}。

\medskip
\textbf{\textcolor{CProof}{方法一：最快证明——两个平方一减，立刻得到半角}}

\textbf{切入点：} 焦距 $2c$ 和双曲线定义给出的两腰之差 $2a$，正好都可以平方。

\textbf{第1步：} 由余弦定理及双曲线定义，直接整体相减：
\[
\begin{aligned}
4b^2
&=4c^2-4a^2\\
&=(m^2+n^2-2mn\cos\theta)-(m-n)^2\\
&=2mn(1-\cos\theta)\\
&=4mn\sin^2\frac{\theta}{2}.
\end{aligned}
\]
于是
\[\boxed{b^2=mn\sin^2\frac{\theta}{2}}.\]
\textbf{第2步：} 面积公式配上 $\sin\theta=2\sin(\theta/2)\cos(\theta/2)$：
\[
\begin{aligned}
S&=\tfrac12 mn\sin\theta
=mn\sin\tfrac\theta2\cos\tfrac\theta2\\
&=\underbrace{mn\sin^2\tfrac\theta2}_{b^2}
  \frac{\cos(\theta/2)}{\sin(\theta/2)}
=\boxed{b^2\cot\frac\theta2}.
\end{aligned}
\]
\textbf{速记：$4c^2-4a^2=4b^2$；余弦定理一减；面积配半角。}

\medskip
\end{proofbox}
\clearpage
\begin{proofbox}
\textbf{\textcolor{CProof}{方法二：最低门槛——垂线把焦点三角形变成勾股问题}}
\begin{center}\begin{tikzpicture}[font=\small,>=stealth]
\draw[black!30,rounded corners=2pt] (0,0) rectangle (14.4,9.2);
\node[anchor=west,font=\small\bfseries] at (.5,8.7) {向焦半径所在直线作垂线：钝角时，垂足落在反向延长线上};
\node at (3.6,7.7) {锐角：$P=(2,3)$};
\node at (10.8,7.7) {钝角：$P=(2/\sqrt3,1)$};
% Local orthonormal frames: P=(0,0), PF1 is positive horizontal.
\begin{scope}[shift={(1.5,4.3)},scale=.75]
\fill[orange!12] (0,0)--(5,0)--(1.8,2.4)--cycle;
\draw[orange!85!black,very thick] (0,0)--(5,0)--(1.8,2.4)--cycle;
\draw[violet,thick] (1.8,0)--(1.8,2.4);
\draw (1.8,.25)--(2.05,.25)--(2.05,0);
\fill (0,0) circle (2.5pt) node[below left=6pt] {$P$};
\fill (5,0) circle (2.5pt) node[right=6pt] {$F_1$};
\fill (1.8,2.4) circle (2.5pt) node[above=8pt] {$F_2$};
\fill (1.8,0) circle (2.5pt) node[below right=6pt] {$H$};
\node[left=7pt] at (.9,1.2) {$n$};
\node[right=6pt] at (1.8,1.2) {$v$};
\node[above right=5pt] at (3.4,1.2) {$2c=4$};
\node at (.9,-1.4) {$u>0$};
\node at (3.4,-1.4) {$m-u$};
\draw (.7,0) arc[start angle=0,end angle=53.130102,radius=.7];
\node at (.85,.45) {$\theta$};
\end{scope}
\pgfmathsetmacro{\nn}{4/sqrt(3)-1}
\pgfmathsetmacro{\mm}{4/sqrt(3)+1}
\pgfmathsetmacro{\uu}{-5*\nn/13}
\pgfmathsetmacro{\vv}{12*\nn/13}
\begin{scope}[shift={(9.8,4.3)},scale=.75]
\draw[black!45,dashed] (\uu,0)--(0,0);
\fill[orange!12] (0,0)--(\mm,0)--(\uu,\vv)--cycle;
\draw[orange!85!black,very thick] (0,0)--(\mm,0)--(\uu,\vv)--cycle;
\draw[violet,thick] (\uu,0)--(\uu,\vv);
\draw (\uu,.2)--({\uu+.2},.2)--({\uu+.2},0);
\fill (0,0) circle (2.5pt) node[below right=6pt] {$P$};
\fill (\mm,0) circle (2.5pt) node[right=6pt] {$F_1$};
\fill (\uu,\vv) circle (2.5pt) node[above left=8pt] {$F_2$};
\fill (\uu,0) circle (2.5pt) node[below left=6pt] {$H$};
\node[left=6pt] at (\uu,{\vv/2}) {$v$};
\node[above right=5pt] at ({(\mm+\uu)/2},{\vv/2}) {$2c=4$};
\node at (-.7,-1.4) {$u<0$};
\node at (2.4,-1.4) {$HF_1=m-u$};
\draw (.6,0) arc[start angle=0,end angle=112.619865,radius=.6];
\node at (.15,.35) {$\theta$};
\end{scope}
\node at (3.6,2.5) {$0<u<m$，$H$ 在 $PF_1$ 上};
\node at (10.8,2.5) {$u<0$，$H$ 在 $P$ 的左侧};
\node at (7.2,1.5) {$u=n\cos\theta$ 是有向投影，$v=n\sin\theta>0$ 是垂距；两图保持同尺度。};
\node at (7.2,.6) {勾股统一写成：$(2c)^2=(m-u)^2+v^2=(m-n\cos\theta)^2+n^2\sin^2\theta$。};
\end{tikzpicture}\end{center}

先把较远的焦点记为 $F_1$，使 $m>n$；交换焦点名称不改变面积或顶角。
图中只把三角形平移、旋转（必要时反射），不改变长度或角度。
记有向投影 $u=n\cos\theta$，垂距 $v=n\sin\theta>0$。
无论 $H$ 在线段上还是延长线上，$HF_1=m-u>0$，勾股均给出
\[
 (2c)^2=(m-u)^2+v^2=m^2+n^2-2mn\cos\theta.
\]
减去双曲线定义给出的 $(m-n)^2=4a^2$：
\[
 4b^2=(2c)^2-(m-n)^2
 =2mn(1-\cos\theta)=4mn\sin^2\frac\theta2.
\]
以 $PF_1$ 为底，图中的紫色垂距为高，故
\[
 S=\frac12 m\,v=\frac12mn\sin\theta
 =\boxed{b^2\cot\frac\theta2}.
\]
\textbf{读图提醒：}$u$ 可以为负，$v$ 才是始终为正的高度；不能把投影当作无符号长度。
\end{proofbox}
\clearpage
\begin{proofbox}
\textbf{\textcolor{CProof}{方法三：常规解析几何——设 $P(x,y)$，求半角}}

先考虑横向标准式
$\frac{x^2}{a^2}-\frac{y^2}{b^2}=1$，通过轴对称性只需讨论右支上方的 $P(x,y)$（$x>a,y>0$），焦点为 $F_1(-c,0)$、$F_2(c,0)$。

\textbf{第1步（底高）：}
\[S=\frac12(2c)y=cy.\]
\textbf{第2步（两条焦半径的和）：} 记 $m=PF_1$、$n=PF_2$，在右支上 $m-n=2a$；又因为
\[
m^2-n^2=((x+c)^2+y^2)-((x-c)^2+y^2)=4cx,
\]
所以
\[
(m-n)(m+n)=4cx
\quad\Longrightarrow\quad
\boxed{m+n=\frac{2cx}{a}}.
\]
\textbf{第3步（把半角正切变成边长）：} 由余弦定理与面积公式分别得到
\[
1+\cos\theta=\frac{(m+n)^2-(2c)^2}{2mn},
\qquad \sin\theta=\frac{2S}{mn}.
\]
因此
\[
\tan\frac\theta2
=\frac{\sin\theta}{1+\cos\theta}
=\frac{4S}{(m+n)^2-4c^2}.
\]
\textbf{第4步（用双曲线方程消 $x$）：}
\[
(m+n)^2-4c^2
=4c^2\left(\frac{x^2}{a^2}-1\right)
=\frac{4c^2y^2}{b^2}.
\]
于是
\[
\tan\frac\theta2
=\frac{4cy}{4c^2y^2/b^2}
=\frac{b^2}{cy}
=\frac{b^2}{S}.
\]
整理即 $\boxed{S=b^2\cot(\theta/2)}$。
\textbf{对称推广：} 另一支或 $y<0$ 时用 $|y|$；纵向标准式交换 $x,y$ 即可。

\medskip
\end{proofbox}
\clearpage
\begin{proofbox}
\textbf{\textcolor{CProof}{方法四：常规三角形法——完整写出联立消元}}

这是解答题可直接写出的六行规范证明。
\[
\begin{aligned}
&m^2+n^2-2mn=(m-n)^2=4a^2,\\
&m^2+n^2-2mn\cos\theta=4c^2,\\
&2mn(1-\cos\theta)=4(c^2-a^2)=4b^2,\\
&mn=\frac{2b^2}{1-\cos\theta},\\
&S=\frac12 mn\sin\theta
=\frac{b^2\sin\theta}{1-\cos\theta},\\
&\boxed{S=b^2\cot\frac\theta2}.
\end{aligned}
\]
其中最后一步只需记住
\[
\frac{\sin\theta}{1-\cos\theta}
=\frac{2\sin(\theta/2)\cos(\theta/2)}{2\sin^2(\theta/2)}
=\cot\frac\theta2.
\]
\textbf{与方法一的区别：} 方法一直接得到 $b^2=mn\sin^2(\theta/2)$，心算快；方法四先求 $mn$，便于按教材规范书写和检查。

\medskip
\textbf{\textcolor{CProof}{对照推导：为什么椭圆是 $\tan$，双曲线是 $\cot$？}}
任意焦点三角形都有
\[
(2c)^2=(m-n)^2+4mn\sin^2\frac\theta2
=(m+n)^2-4mn\cos^2\frac\theta2.
\]
\begin{itemize}[leftmargin=*,itemsep=0.3em]
\item \textbf{椭圆：} $m+n=2a$、$b^2=a^2-c^2$，由右式得 $b^2=mn\cos^2(\theta/2)$；因此 $S=b^2\tan(\theta/2)$。
\item \textbf{双曲线：} $|m-n|=2a$、$b^2=c^2-a^2$，由左式得 $b^2=mn\sin^2(\theta/2)$；因此 $S=b^2\cot(\theta/2)$。
\end{itemize}
\textbf{两条面积公式的区别，根本上来自“和固定”与“差固定”。}
\end{proofbox}

\clearpage
\begin{proofbox}
\textbf{\textcolor{CProof}{方法五：纯几何法——先认清三个旁切圆（不使用余弦定理）}}

把离 $P$ 较远的焦点记为 $F_1$，较近的记为 $F_2$。设
\[
PF_1=m,\quad PF_2=n,\quad m-n=2a,\quad F_1F_2=2c,
\qquad s=\frac{m+n+2c}{2}.
\]
\textbf{先补充旁切圆性质。} 每个非退化三角形\textbf{恰有三个旁切圆}，分别对应三个顶点。
每个圆与该顶点的\textbf{对边线段}及另外两边的\textbf{延长线}相切。
圆心位于该顶点的内角平分线及另两个顶点的外角平分线上；到三边所在直线的距离都等于半径。
同一点向同一圆所引的两条切线段等长，且半径垂直于切线。

\textbf{（1） 本法选 $F_1$ 对应的旁切圆。} 记圆心 $I_1$、半径 $r_1$。
\begin{center}
\begin{tikzpicture}[x=1cm,y=1cm,line cap=round,line join=round,font=\small]
\draw[black!25,rounded corners=3pt] (0,0) rectangle (14.4,10.5);
\node[font=\small\bfseries] at (7.2,10.0) {（1） 对应 $F_1$：与边 $PF_2$ 相切（本法选它）};
\begin{scope}[shift={(3.00000000,2.20000000)},x=1.02cm,y=1.02cm]
\coordinate (A) at (-2.00000000,0.00000000);
\coordinate (B) at (2.50000000,3.96862697);
\coordinate (C) at (2.00000000,0.00000000);
\coordinate (I) at (5.00000000,2.64575131);
\coordinate (T) at (2.37500000,2.97647022);
\coordinate (U) at (3.25000000,4.63006479);
\coordinate (V) at (5.00000000,0.00000000);
\fill[green!4] (I) circle (2.64575131);
\draw[green!45!black,thick] (I) circle (2.64575131);
\fill[red!7] (A)--(B)--(C)--cycle;
\draw[red!75!black,thick] (A)--(B)--(C)--cycle;
\draw[red!65!black,dashed] (B)--(U);
\draw[red!65!black,dashed] (C)--(V);
\draw[purple!75,dashed] (I)--(T);
\draw[purple!75] (2.39093137,3.10292157)--(2.51738272,3.08699020)--(2.50145135,2.96053885);
\draw[purple!75,dashed] (I)--(U);
\draw[purple!75] (3.34558824,4.71436569)--(3.42988913,4.61877746)--(3.33430090,4.53447656);
\draw[purple!75,dashed] (I)--(V);
\draw[purple!75] (4.87254902,0.00000000)--(4.87254902,0.12745098)--(5.00000000,0.12745098);
\fill (A) circle (1.5pt); \node[below left] at (A) {$F_1$};
\fill (B) circle (1.5pt); \node[above left] at (B) {$P$};
\fill (C) circle (1.5pt); \node[below right] at (C) {$F_2$};
\fill (I) circle (1.5pt); \node[above right] at (I) {$I_1$};
\fill[purple!75] (T) circle (1.5pt); \node[left,fill=white,inner sep=1pt] at (T) {$T$};
\fill[purple!75] (U) circle (1.5pt); \node[above left,fill=white,inner sep=1pt] at (U) {$U$};
\fill[purple!75] (V) circle (1.5pt); \node[below,fill=white,inner sep=1pt] at (V) {$V$};
\end{scope}
\node at (7.2,1.15) {$T\in\overline{PF_2}$；$U$ 越过 $P$；$V$ 越过 $F_2$};
\node at (7.2,.55) {$r_1=\sqrt7,\quad PT=PU=1=c-a$};
\end{tikzpicture}
\end{center}

后文简记 $I=I_1,\ r=r_1$，不与另外两个圆的半径混用。
\end{proofbox}

\clearpage
\begin{proofbox}
\textbf{\textcolor{CProof}{方法五（预备知识续）：另外两个旁切圆}}

\textbf{（2） $P$ 对应的旁切圆：} 与边 $F_1F_2$ 相切，圆心记为 $I_P$，半径记为 $r_P$。
\begin{center}
\begin{tikzpicture}[x=1cm,y=1cm,line cap=round,line join=round,font=\small]
\draw[black!25,rounded corners=3pt] (0,0) rectangle (14.4,10.5);
\node[font=\small\bfseries] at (7.2,10.0) {（2） 对应 $P$：与边 $F_1F_2$ 相切};
\begin{scope}[shift={(7.00000000,6.00000000)},x=0.8cm,y=0.8cm]
\coordinate (A) at (-2.00000000,0.00000000);
\coordinate (B) at (2.50000000,3.96862697);
\coordinate (C) at (2.00000000,0.00000000);
\coordinate (I) at (-1.00000000,-2.64575131);
\coordinate (T) at (-1.00000000,0.00000000);
\coordinate (U) at (-2.75000000,-0.66143783);
\coordinate (V) at (1.62500000,-2.97647022);
\fill[green!4] (I) circle (2.64575131);
\draw[green!45!black,thick] (I) circle (2.64575131);
\fill[red!7] (A)--(B)--(C)--cycle;
\draw[red!75!black,thick] (A)--(B)--(C)--cycle;
\draw[red!65!black,dashed] (A)--(U);
\draw[red!65!black,dashed] (C)--(V);
\draw[purple!75,dashed] (I)--(T);
\draw[purple!75] (-0.86315789,0.00000000)--(-0.86315789,-0.13684211)--(-1.00000000,-0.13684211);
\draw[purple!75,dashed] (I)--(U);
\draw[purple!75] (-2.64736842,-0.57092528)--(-2.55685588,-0.67355686)--(-2.65948746,-0.76406941);
\draw[purple!75,dashed] (I)--(V);
\draw[purple!75] (1.60789474,-3.11223904)--(1.47212592,-3.09513378)--(1.48923118,-2.95936496);
\fill (A) circle (1.5pt); \node[above left] at (A) {$F_1$};
\fill (B) circle (1.5pt); \node[above left] at (B) {$P$};
\fill (C) circle (1.5pt); \node[below right] at (C) {$F_2$};
\fill (I) circle (1.5pt); \node[above right] at (I) {$I_P$};
\fill[purple!75] (T) circle (1.5pt); \node[left,fill=white,inner sep=1pt] at (T) {$T$};
\fill[purple!75] (U) circle (1.5pt); \node[above left,fill=white,inner sep=1pt] at (U) {$U$};
\fill[purple!75] (V) circle (1.5pt); \node[below,fill=white,inner sep=1pt] at (V) {$V$};
\end{scope}
\node at (7.2,1.15) {$T\in\overline{F_1F_2}$；$U$ 越过 $F_1$；$V$ 越过 $F_2$};
\node at (7.2,.55) {$r_P=\sqrt7$（本例等腰，恰有 $r_P=r_1$）};
\end{tikzpicture}
\end{center}

它的圆心在 $P$ 的\textbf{内角平分线}上；因此不能把本法使用的外半角关系直接套给它。

\textbf{三图均用同一个边长为 $6,4,4$ 的三角形。} 三图缩放不同，各自均等比例绘制；
比较半径请看数值，不按印在纸上的大小比较。本例等腰，两个旁切圆半径相等，但仍是两个不同的圆。
\end{proofbox}

\clearpage
\begin{proofbox}
\textbf{\textcolor{CProof}{方法五（预备知识续）：第三个旁切圆与面积关系}}

\textbf{（3） $F_2$ 对应的旁切圆：} 与边 $F_1P$ 相切，圆心记为 $I_2$，半径记为 $r_2$。
\begin{center}
\begin{tikzpicture}[x=1cm,y=1cm,line cap=round,line join=round,font=\small]
\draw[black!25,rounded corners=3pt] (0,0) rectangle (14.4,10.5);
\node[font=\small\bfseries] at (7.2,10.0) {（3） 对应 $F_2$：与边 $F_1P$ 相切};
\begin{scope}[shift={(8.60000000,2.05000000)},x=0.46cm,y=0.46cm]
\coordinate (A) at (-2.00000000,0.00000000);
\coordinate (B) at (2.50000000,3.96862697);
\coordinate (C) at (2.00000000,0.00000000);
\coordinate (I) at (-5.00000000,7.93725393);
\coordinate (T) at (0.25000000,1.98431348);
\coordinate (U) at (2.87500000,6.94509719);
\coordinate (V) at (-5.00000000,0.00000000);
\fill[green!4] (I) circle (7.93725393);
\draw[green!45!black,thick] (I) circle (7.93725393);
\fill[red!7] (A)--(B)--(C)--cycle;
\draw[red!75!black,thick] (A)--(B)--(C)--cycle;
\draw[red!65!black,dashed] (B)--(U);
\draw[red!65!black,dashed] (A)--(V);
\draw[purple!75,dashed] (I)--(T);
\draw[purple!75] (0.03804348,1.79738540)--(-0.14888460,2.00934192)--(0.06307192,2.19627001);
\draw[purple!75,dashed] (I)--(U);
\draw[purple!75] (2.83967391,6.66470507)--(2.55928179,6.70003116)--(2.59460788,6.98042328);
\draw[purple!75,dashed] (I)--(V);
\draw[purple!75] (-5.28260870,0.00000000)--(-5.28260870,0.28260870)--(-5.00000000,0.28260870);
\fill (A) circle (1.5pt); \node[below left] at (A) {$F_1$};
\fill (B) circle (1.5pt); \node[right] at (B) {$P$};
\fill (C) circle (1.5pt); \node[below right] at (C) {$F_2$};
\fill (I) circle (1.5pt); \node[above right] at (I) {$I_2$};
\fill[purple!75] (T) circle (1.5pt); \node[left,fill=white,inner sep=1pt] at (T) {$T$};
\fill[purple!75] (U) circle (1.5pt); \node[above left,fill=white,inner sep=1pt] at (U) {$U$};
\fill[purple!75] (V) circle (1.5pt); \node[below,fill=white,inner sep=1pt] at (V) {$V$};
\end{scope}
\node at (7.2,1.15) {$T\in\overline{F_1P}$；$U$ 越过 $P$；$V$ 越过 $F_1$};
\node at (7.2,.55) {$r_2=3\sqrt7$（大圆仍完整画出）};
\end{tikzpicture}
\end{center}

\textbf{统一记法：} 对边长为 $d$ 的旁切圆半径 $\rho$ 满足 $S=(s-d)\rho$。
这是把三边作底、半径作高，两个三角形面积相加再减一个得到的（后面画出本法所用的分解）。因此
\[
\begin{array}{c|c|c}
\text{对应顶点}&\text{实际相切的边}&\text{面积关系}\\ \hline
F_1&PF_2=n&S=(s-n)r_1=(a+c)r_1\\
P&F_1F_2=2c&S=(s-2c)r_P\\
F_2&PF_1=m&S=(s-m)r_2=(c-a)r_2
\end{array}
\]
\textbf{选择理由：} 本法用 $I_1$，是为了让切线长得到 $c-a$、面积因子得到 $c+a$，两者相乘恰是 $b^2$。
\end{proofbox}

\clearpage
\begin{proofbox}
\textbf{\textcolor{CProof}{方法五（正式推导）：等切线先给出 $c-a$}}

\textbf{第1步：回到第（1）图，只用同点切线等长。}
\[
PT=PU=x,\quad F_2T=F_2V=n-x,\quad F_1U=F_1V.
\]
由延长线上的线段相加，
\[
\underbrace{m+x}_{F_1U}=\underbrace{2c+n-x}_{F_1V}
\quad\Longrightarrow\quad
\boxed{PT=x=c-a}.
\]
\textbf{第2步：外角平分线给出半径。} $PU$ 与 $PF_1$ 反向，故外角为 $\pi-\theta$。
\begin{center}
\begin{tikzpicture}[x=1cm,y=1cm,line cap=round,line join=round,font=\small]
\draw[black!25,rounded corners=3pt] (0,0) rectangle (14.4,9.8);
\node[font=\small\bfseries] at (7.2,9.3) {外角半角局部图：$PI$ 平分外角，不是内角};
\begin{scope}[shift={(4,3)},x=1.65cm,y=1.65cm]
\coordinate (P) at (0,0); \coordinate (T) at (1,0);
\coordinate (I) at (1,2.64575131);
\coordinate (U) at (-.75,0.66143783);
\fill[green!6] (P)--(T)--(I)--cycle;
\draw[red!75!black,thick] (P)--(1.7,0) (P)--(U);
\draw[black!50,->] (P)--(-41.409622:1.5) node[right] {指向 $F_1$};
\draw[green!45!black,thick,dashed] (P)--(I);
\draw[purple!75,thick] (I)--(T);
\draw[purple!75] (.86,0)--(.86,.14)--(1,.14);
\draw[green!45!black] (.52,0) arc(0:69.295189:.52);
\draw[green!45!black] (69.295189:.69) arc(69.295189:138.590378:.69);
\node at (.67,.45) {$\alpha$}; \node at (-.25,.88) {$\alpha$};
\draw[red!75!black] (-41.409622:.83) arc(-41.409622:0:.83);
\node at (1.02,-.42) {$\theta$};
\node[below left] at (P) {$P$}; \node[below right] at (T) {$T$};
\node[above left] at (I) {$I=I_1$}; \node[above left] at (U) {$U$};
\node[right] at (1,1.6) {$r$};
\node at (-.8,-.6) {$PT=c-a$};
\end{scope}
\node[align=left,anchor=north west] at (9,6.7) {$2\alpha+\theta=\pi$\\[1em]$\alpha=\dfrac{\pi-\theta}{2}$\\[1em]$\dfrac r{PT}=\tan\alpha$\\[1em]$\phantom{\dfrac r{PT}}=\cot\dfrac\theta2$};
\node at (7.2,.55) {只保留原图的局部关系，经旋转、放大；$PU$ 与 $PF_1$ 方向相反。};
\end{tikzpicture}
\end{center}

在直角三角形 $PIT$ 中，$\angle IPT=(\pi-\theta)/2$，所以
\[
\frac r{PT}=\tan\frac{\pi-\theta}{2}=\cot\frac\theta2,
\qquad \boxed{r=(c-a)\cot\frac\theta2}.
\]
这里用的是\textbf{外角的一半}，不是 $\theta/2$。
\end{proofbox}

\clearpage
\begin{proofbox}
\textbf{\textcolor{CProof}{方法五（完成）：面积“加、加、减”给出 $c+a$}}

\textbf{第3步：同一个四边形，两种分法。} 下面两图的四边形 $PF_1F_2I$ 完全相同、尺度相同，只更换分割线。
\begin{center}
\begin{tikzpicture}[x=1cm,y=1cm,line cap=round,line join=round,font=\small]
\draw[black!25,rounded corners=3pt] (0,0) rectangle (14.4,6.8);
\node[font=\small\bfseries] at (7.2,6.3) {第一种分法：沿 $F_1I$ 分成蓝、绿两个三角形};
\begin{scope}[shift={(4.5,1.35)},x=1.05cm,y=1.05cm]
\coordinate (A) at (-2,0);
\coordinate (P) at (2.5,{3*sqrt(7)/2});
\coordinate (B) at (2,0);
\coordinate (I) at (5,{sqrt(7)});
\fill[blue!23] (P)--(A)--(I)--cycle;
\fill[green!23] (I)--(A)--(B)--cycle;
\draw[blue!65!black,line width=1.2pt] (A)--(I);
\node[blue!60!black] at (1.9,2.5) {$S_{PF_1I}$};
\node[green!40!black] at (1.7,.73) {$S_{IF_1F_2}$};
\draw[black!80,line width=1.1pt] (P)--(A)--(B)--(I)--cycle;
\foreach \q in {A,P,B,I} \fill (\q) circle (1.7pt);
\node[below left] at (A) {$F_1$};
\node[above left] at (P) {$P$};
\node[below] at (B) {$F_2$};
\node[right] at (I) {$I$};
\end{scope}
\node at (7.2,.5) {$S_{\text{四边形}}=\textcolor{blue!60!black}{S_{PF_1I}}+\textcolor{green!40!black}{S_{IF_1F_2}}$};
\end{tikzpicture}
\end{center}

\begin{center}
\begin{tikzpicture}[x=1cm,y=1cm,line cap=round,line join=round,font=\small]
\draw[black!25,rounded corners=3pt] (0,0) rectangle (14.4,6.8);
\node[font=\small\bfseries] at (7.2,6.3) {第二种分法：沿 $PF_2$ 分成橙、红两个三角形};
\begin{scope}[shift={(4.5,1.35)},x=1.05cm,y=1.05cm]
\coordinate (A) at (-2,0);
\coordinate (P) at (2.5,{3*sqrt(7)/2});
\coordinate (B) at (2,0);
\coordinate (I) at (5,{sqrt(7)});
\fill[orange!23] (P)--(A)--(B)--cycle;
\fill[red!23] (I)--(P)--(B)--cycle;
\draw[red!70!black,line width=1.2pt] (P)--(B);
\node[orange!60!black] at (.55,1.25) {$S_{PF_1F_2}$};
\node[red!60!black] at (3.3,2.25) {$S_{IPF_2}$};
\draw[black!80,line width=1.1pt] (P)--(A)--(B)--(I)--cycle;
\foreach \q in {A,P,B,I} \fill (\q) circle (1.7pt);
\node[below left] at (A) {$F_1$};
\node[above left] at (P) {$P$};
\node[below] at (B) {$F_2$};
\node[right] at (I) {$I$};
\end{scope}
\node at (7.2,.5) {$S_{\text{四边形}}=\textcolor{orange!60!black}{S_{PF_1F_2}}+\textcolor{red!60!black}{S_{IPF_2}}$};
\end{tikzpicture}
\end{center}

两种分法的总面积相等，故\textbf{把第二图的红色三角形减去}：
\[
S_{PF_1F_2}=S_{PF_1I}+S_{IF_1F_2}-S_{IPF_2}.
\]
\end{proofbox}
\clearpage
\begin{proofbox}
\textbf{\textcolor{CProof}{方法五（最后一步）：将面积分解换成底乘高}}

三个底边所在直线到 $I$ 的距离都为 $r$，因此
\[
S=\frac12 mr+\frac12(2c)r-\frac12 nr
 =\frac r2(m-n+2c)=\boxed{(a+c)r}.
\]
\textbf{第4步：代入半径，认出 $b^2$。}
\[
S=(a+c)(c-a)\cot\frac\theta2
=(c^2-a^2)\cot\frac\theta2
=\boxed{b^2\cot\frac\theta2}.
\]
\textbf{记忆：} 等切线找 $c-a$，面积分解找 $c+a$，乘起来就是 $b^2$。

\textbf{适用范围：} $P$ 在任一分支且不是实轴顶点时均成立；只需先按距离远近重标 $F_1,F_2$。
顶点处三角形退化，本旁切圆构造不适用；面积公式仍可按极限延拓为 $0$。
\end{proofbox}
